\documentclass[11pt]{article}
\usepackage{amsmath,amssymb,enumerate,xcolor}
\usepackage[margin=1.2in]{geometry}
\usepackage{fancyhdr}
\pagestyle{fancy}
\setlength{\headheight}{20pt}
\setlength{\headsep}{20pt}
\fancyhead[L]{\small CS 171, Fall 2026}
\fancyhead[R]{\small Prof. Sanjam Garg}
\begin{document}
\section*{CS 171: Problem Set 8}

\section{Bit commitments from a PRG (15 Points)}
Let $G:\{0,1\}^{n}\to\{0,1\}^{3n}$ be a secure pseudorandom generator. Consider the following protocol:
\begin{enumerate}
\item The receiver samples $R\leftarrow\{0,1\}^{3n}$ uniformly and sends $R$ to the sender.
\item To commit to $b\in\{0,1\}$, the sender samples $s\leftarrow\{0,1\}^{n}$ uniformly and sends
\[c=G(s)\oplus bR,\]
where $0R=0^{3n}$ and $1R=R$.
\item To open, the sender reveals $(b,s)$. The receiver accepts exactly when $c=G(s)\oplus bR$.
\end{enumerate}
\begin{enumerate}[(a)]
\item Prove correctness.
\item Prove computational hiding, even against a malicious receiver that chooses $R$ arbitrarily before receiving $c$. Give a reduction to PRG security.
\item Prove statistical binding: with probability at least $1-2^{-n}$ over an honest receiver's choice of $R$, no commitment string admits valid openings to both bits---even for a computationally unbounded sender.
\end{enumerate}

\section{Reusing randomness in Pedersen commitments (15 Points)}
Let $\mathbb G$ be a cyclic group of prime order $q>2$, with public generators $g,h$, where computing $\log_g h$ is hard.

Alice wants to commit to two bits $b_1,b_2\in\{0,1\}$. She samples a single $r\leftarrow\mathbb Z_q$ uniformly and publishes
\[C_1=g^{b_1}h^r,\qquad C_2=g^{b_2}h^r.\]
To open, she reveals $(b_1,b_2,r)$, and the receiver checks both equations.
\begin{enumerate}[(a)]
\item Prove that, viewed separately, each $C_i$ perfectly hides its corresponding bit. Does this alone imply that the pair hides both bits?
\item Compute $C_1C_2^{-1}$. Show how a receiver can determine whether the bits are equal and, when they differ, recover both bits.
\item Does this construction remain computationally binding to the ordered pair $(b_1,b_2)$? Prove your answer using the binding of ordinary Pedersen commitments.
\item Alice proposes using $r$ for the first commitment and $r+1\pmod q$ for the second. Does this repair hiding? Justify your answer.
\end{enumerate}

\section{Finding Relations and Discrete Logarithms (10 Points)}
Let $\mathbb G$ be a cyclic group of prime order $q$, generated by $g$, and let $n$ be polynomially bounded in the security parameter. A PPT adversary $\mathcal A$ receives independent uniformly random elements $g_1,\ldots,g_n\in\mathbb G$. It succeeds if it outputs a nonzero vector $(\alpha_1,\ldots,\alpha_n)\in\mathbb Z_q^n$ such that
\[\prod_{i=1}^n g_i^{\alpha_i}=1.\]

Suppose $\mathcal A$ succeeds with non-negligible probability. Prove that there is a PPT algorithm $\mathcal B$ that, given $h=g^x$ for uniformly random $x\in\mathbb Z_q$, computes $x$ with non-negligible probability as well. More precisely, if $\mathcal A$ succeeds with probability $\epsilon$, show that $\mathcal B$ succeeds with probability at least $\epsilon-1/q$. Conclude that finding such a relation is hard under the discrete-log assumption.

\textit{Hint:} Recall the binding proof for Pedersen commitments.


\end{document}
